% !TEX root = ../tree_notes.tex
% Roadmap: parts at a glance (numbers are automatic), dependency map, study plan.
\chapter{Roadmap}

\section*{The book at a glance}

The book has five parts, ordered by difficulty. The chapter and problem counts
below are counted automatically from the chapters themselves; an outline
chapter contributes its planned problems until it is written.

\begin{center}
\renewcommand{\arraystretch}{1.25}
\begin{tabular}{@{}llllr@{}}
\toprule
\textbf{Part} & \textbf{Title} & \textbf{Level} & \textbf{Chapters} & \textbf{Problems} \\
\midrule
I   & Foundations               & Beginner -- LC Easy/Medium & 1--6   & \stat{partI:problems} \\
II  & Ordered Trees, Heaps, Tries & LC Medium -- CF 1600     & 7--11  & \stat{partII:problems} \\
III & Algorithms on Rooted Trees & CF 1400--2000             & 12--17 & \stat{partIII:problems} \\
IV  & Query Trees               & CF 1600--2300             & 18--23 & \stat{partIV:problems} \\
V   & Expert Techniques         & CF 1800--2800             & 24--31 & \stat{partV:problems} \\
\midrule
    & \textbf{Total}          &                     & \textbf{31} & \textbf{\stat{book:problems}} \\
\bottomrule
\end{tabular}
\end{center}

\textbf{Progress.} \stat{book:writtenchapters} of 31 chapters are written,
with \stat{book:written} problems on paper (\stat{book:solved} of them fully
solved) and \stat{book:planned} more planned in the outline.

Part~I builds the vocabulary every later chapter uses: what a tree is, how
binary trees are shaped, how to walk them, and how to store them. Read it
first even if you already ``know'' trees -- the notation is fixed here.

\section*{How the parts depend on each other}

\begin{figure}[H]
\centering
\begin{tikzpicture}[diagram, node distance=5mm and 5mm,
    part/.style={boxnode, text width=25mm, minimum height=12mm, fill=primaryblue!10, draw=primaryblue},
    keys/.style={annot, text width=26mm, anchor=north}]
  \node[part] (p1) {\textbf{I} Foundations\\\footnotesize Ch.\ 1--6};
  \node[part, right=of p1] (p2) {\textbf{II} Ordered \& tries\\\footnotesize Ch.\ 7--11};
  \node[part, right=of p2] (p3) {\textbf{III} Rooted algos\\\footnotesize Ch.\ 12--17};
  \node[part, right=of p3] (p4) {\textbf{IV} Query trees\\\footnotesize Ch.\ 18--23};
  \node[part, right=of p4] (p5) {\textbf{V} Expert\\\footnotesize Ch.\ 24--31};
  \foreach \a/\b in {p1/p2,p2/p3,p3/p4,p4/p5} \draw[flow] (\a) -- (\b);
  \node[keys] at ([yshift=-2mm]p1.south) {terminology, recursion, traversals, storage, Euler tour, construction};
  \node[keys] at ([yshift=-2mm]p2.south) {BST, balanced, Cartesian, heaps, tries};
  \node[keys] at ([yshift=-2mm]p3.south) {diameter, tree DP, rerooting, LCA, DSU/sack};
  \node[keys] at ([yshift=-2mm]p4.south) {Fenwick, segment, lazy, persistent, HLD, Li Chao};
  \node[keys] at ([yshift=-2mm]p5.south) {centroid, dynamic trees, virtual, offline, isomorphism, graph-trees, Prufer, MST};
  \draw[->, very thick, accentorange] ([yshift=-19mm]p1.south west) -- ([yshift=-19mm]p5.south east)
      node[annot, pos=0.5, below=1mm] {difficulty grows from LC Easy to CF 2800+};
\end{tikzpicture}
\caption*{Each part assumes the previous ones. Under each part: the areas it covers.}
\end{figure}

\section*{Study plan}

\begin{tcolorbox}[enhanced, breakable, colback=lightgray, colframe=primaryblue, title={Beginner track -- Chapters 1--6 (3--4 weeks)}]
\textbf{Focus.} Tree vocabulary; binary tree shapes and node counts; DFS and BFS traversals (recursive and iterative); building trees from arrays and traversals; flattening a subtree to an interval with an Euler tour; rebuilding a tree from parent arrays and degrees.\par
\textbf{Goals.} Traverse any tree four ways without thinking; compute depth, height and size by recursion; solve every $\lv1$--$\lv3$ problem (LeetCode Easy--Medium) in Chapters 1--6 unaided.
\end{tcolorbox}

\begin{tcolorbox}[enhanced, breakable, colback=lightgray, colframe=successgreen, title={Intermediate track -- Chapters 7--11 (3--4 weeks)}]
\textbf{Focus.} BST search, insert, delete; balanced trees (AVL, red--black ideas); Cartesian trees and the monotonic stack; heaps and priority queues; tries for strings.\par
\textbf{Goals.} Implement a BST and a heap from scratch; know why balance keeps operations $O(\log n)$; reach LC Medium.
\end{tcolorbox}

\begin{tcolorbox}[enhanced, breakable, colback=lightgray, colframe=accentorange, title={Advanced track -- Chapters 12--17 (4--5 weeks)}]
\textbf{Focus.} Diameter and centroids; tree DP (subtree and knapsack forms); rerooting; lowest common ancestor with binary lifting; DSU, small-to-large merging and sack.\par
\textbf{Goals.} Solve rooted-tree problems with a single DFS; answer LCA in $O(\log n)$; reach CF 1400--2000.
\end{tcolorbox}

\begin{tcolorbox}[enhanced, breakable, colback=lightgray, colframe=softpurple, title={Query-structure track -- Chapters 18--23 (4--5 weeks)}]
\textbf{Focus.} Fenwick trees; segment trees and lazy propagation; merge-sort and persistent trees; heavy--light decomposition; Li Chao trees and segment tree beats.\par
\textbf{Goals.} Build, query and update a segment tree fluently; handle path queries on a tree with HLD; reach CF 1600--2300.
\end{tcolorbox}

\begin{tcolorbox}[enhanced, breakable, colback=lightgray, colframe=deepblue, title={Expert track -- Chapters 24--31 (5--6 weeks)}]
\textbf{Focus.} Centroid decomposition; dynamic trees (Euler-tour and link-cut); virtual trees; offline queries (Mo's on trees, Tarjan LCA); tree hashing and isomorphism; bridge, block-cut and functional graphs; Prufer codes and the Matrix--Tree theorem; minimum spanning trees.\par
\textbf{Goals.} Reach CF 1800--2800+. These chapters assume everything before them.
\end{tcolorbox}
